\chapter{Cyberbullying}

\section*{Definition and Overview}
Cyberbullying is bullying that takes place using digital technology, including social media, messaging apps, gaming platforms, and mobile phones. It involves repeated, intentional harmful behaviour such as sending threatening or humiliating messages, spreading rumours online, sharing embarrassing images, or excluding someone from online groups. Because it happens online, cyberbullying can follow a person everywhere, reach a wide audience quickly, and leave a permanent digital record. It affects children, teenagers, and adults, and it can have serious effects on mental health, self-esteem, and safety.

\section*{Epidemiology}
Cyberbullying is widespread, particularly among school-aged children and young adults. Studies suggest that a significant minority of young people -- commonly around one in five to one in four -- experience cyberbullying at some point. It overlaps with face-to-face bullying, and some people experience both. Girls are more likely to be targeted in some studies, and cyberbullying peaks in early adolescence. In Australia, organisations such as the eSafety Commissioner track and respond to serious cyberbullying, which is among the most common complaints made by young people.

\section*{Aetiology and Pathophysiology}
Cyberbullying arises from the same dynamics as other bullying, amplified by technology.

\begin{itemize}
    \item \textbf{Anonymity:} the ability to hide behind screens emboldens harmful behaviour
    \item \textbf{Wide reach:} posts can be shared rapidly to large audiences
    \item \textbf{24/7 access:} digital devices mean there is no escape from the behaviour
    \item \textbf{Permanence:} content can remain online and resurface later
    \item \textbf{Power imbalance:} as with all bullying, there is an imbalance of power
    \item \textbf{Group dynamics:} bystanders can amplify harm by liking, sharing, or joining in
    \item \textbf{Impact mechanism:} repeated humiliation activates stress responses, harming mood, sleep, and self-esteem
\end{itemize}

\section*{Clinical Features}
The signs of cyberbullying may be noticed by the person, parents, teachers, or friends.

\begin{itemize}
    \item Repeated hurtful messages, comments, or images
    \item Being excluded or ignored in online groups
    \item Fake accounts or identity theft used to harass
    \item Threats, intimidation, or blackmail
    \item Private information or images shared without consent
    \item The person becoming withdrawn, anxious, or tearful
    \item Changes in sleep, appetite, or school performance
    \item Avoiding devices or being secretive about online activity
    \item Declining friendships and social withdrawal
\end{itemize}

\section*{Differential Diagnosis}
Not all conflict online is cyberbullying, and other issues are considered.

\begin{itemize}
    \item A single argument or disagreement, which is not bullying unless repeated and intentional
    \item Online grooming or exploitation, which is a criminal matter
    \item Sexting or sharing images, which needs its own response
    \item Mental health problems, including depression and anxiety, which can coexist or follow cyberbullying
    \item Family or school problems causing distress
    \item Digital addiction or problematic internet use
\end{itemize}

\section*{When to Seek Medical Care}
Anyone being cyberbullied should be supported and, if distress is significant, see a GP or counsellor. Parents should take reports seriously and respond calmly and supportively.

\textbf{Seek urgent help if:}
\begin{itemize}
    \item The person talks about self-harm or suicide
    \item There are threats of violence
    \item There is sexual exploitation or grooming
    \item The person is severely distressed, not sleeping, or not coping with daily life
\end{itemize}

\textbf{Call 000 (or local emergency services) for:} immediate danger of self-harm, violence, or a criminal threat. For crisis support in Australia, call Lifeline on 13 11 14 or Kids Helpline on 1800 55 1800.

\section*{Diagnosis and Investigations}
There is no medical test for cyberbullying; assessment focuses on impact and safety.

\begin{itemize}
    \item \textbf{Conversation:} a safe, non-judgemental discussion of what is happening
    \item \textbf{Evidence gathering:} screenshots, links, and dates of the behaviour
    \item \textbf{Impact assessment:} mood, sleep, school or work function, and any self-harm risk
    \item \textbf{Platform reporting:} reporting content to the social media platform
    \item \textbf{Formal reporting:} reporting to the school, workplace, or the eSafety Commissioner
    \item \textbf{Police involvement:} for threats, harassment, or image-based abuse
    \item \textbf{Mental health review:} by a GP or psychologist where distress is significant
\end{itemize}

\section*{Management}
Management involves the person, parents or carers, schools, and platforms.

\begin{itemize}
    \item \textbf{Stopping the behaviour:} blocking the bully, adjusting privacy settings, and not responding to messages
    \item \textbf{Saving evidence:} screenshots before deleting or blocking
    \item \textbf{Platform reporting:} using the reporting tools of the app or website
    \item \textbf{School involvement:} informing the school, which has a duty to address bullying
    \item \textbf{eSafety Commissioner:} can help remove serious cyberbullying material targeted at young people
    \item \textbf{Police:} for threats, harassment, and image-based abuse
    \item \textbf{Emotional support:} counselling, family support, and helping the person rebuild confidence
    \item \textbf{Digital wellbeing:} healthy limits on device use and support for social connections
\end{itemize}

\section*{Complications}
\begin{itemize}
    \item Depression, anxiety, and low self-esteem
    \item Sleep disturbance and school refusal
    \item Self-harm and suicidal thoughts
    \item Social isolation and loss of friendships
    \item Poor school or work performance
    \item Distress for the whole family
    \item Long-term effects on confidence and trust
    \item The person bullying others, who also needs help
\end{itemize}

\section*{Prognosis and Outcomes}
Most people who experience cyberbullying recover well with support, and effective responses stop the behaviour in most cases. Outcomes are better when the bullying is addressed promptly, evidence is preserved, and the person receives emotional support. When cyberbullying continues untreated, the risk of depression, anxiety, and self-harm rises. Schools, platforms, and the eSafety Commissioner provide practical avenues for resolution, and most young people who are supported through the experience regain their confidence and wellbeing.

\section*{Prevention and Self-Care}
\begin{itemize}
    \item Talk openly with children and young people about their online life
    \item Teach safe online habits: privacy settings, careful sharing, and thinking before posting
    \item Keep screenshots and report rather than retaliate
    \item Block and ignore the bully and tell a trusted adult
    \item Take breaks from devices and maintain offline friendships and activities
    \item Parents: stay involved, use parental controls, and respond supportively to reports
    \item Schools: promote anti-bullying policies and digital citizenship
    \item Seek counselling early if distress is significant
\end{itemize}

\section*{Sources and Further Reading}
The following reputable sources provide further information for healthcare students and patients seeking reliable, current advice about cyberbullying.

\begin{itemize}
    \item eSafety Commissioner. \textit{Cyberbullying help and reporting.} https://www.esafety.gov.au/
    \item Kids Helpline. \textit{Cyberbullying support for young people.} https://kidshelpline.com.au/
    \item Bullying. No Way! \textit{Resources for schools and families.} https://bullyingnoway.gov.au/
    \item Headspace. \textit{Cyberbullying and mental health.} https://headspace.org.au/
    \item ReachOut. \textit{Cyberbullying help for young people.} https://au.reachout.com/
    \item Raising Children Network. \textit{Cyberbullying: how to help.} https://raisingchildren.net.au/
\end{itemize}
